% 本文件由 tools/build_book.py 从 dist/reports/assertions.json 生成，请勿手改。
\begin{longtable}{@{}p{0.06\textwidth}p{0.32\textwidth}p{0.22\textwidth}p{0.16\textwidth}@{}}
\caption{章节 1.4 数值断言（生成时间 2026-10-19T16:08:50-0400）}\\
\toprule
类别 & 断言 & 理论值 & 实测值 \\
\midrule
\endfirsthead
\toprule
类别 & 断言 & 理论值 & 实测值 \\
\midrule
\endhead
\bottomrule
\endlastfoot

A & 幂法实测收敛因子 vs |$\lambda$$_2$/$\lambda$$_1$| & 实测比 =\allowbreak{} r(\allowbreak{}相对误差 < 2\%) & 最大相对误差 0.0000；Rayleigh 商最大误差 0.00e+00；最多 142 步 \checkmark \\
A & 闭式谱对拍(Laplacian/Jordan 型) & max|$\lambda$\_QR $-$ (\allowbreak{}2$-$2cos(\allowbreak{}j$\pi$/\allowbreak{}(\allowbreak{}n+1)))| $\le$ 1e-12 & 最大误差 3.64e-15；剪切矩阵谱 =\allowbreak{} [1. 1. 1.]$\dots$；N =\allowbreak{} A$-$I 的谱 =\allowbreak{} [0. 0. 0.]$\dots$ \checkmark \\
A & Hessenberg 化的结构/正交/谱三项误差 & 下三角(\allowbreak{}$\le$$-$2) =\allowbreak{} 0、\allowbreak{}$\|$Q$^{\mathsf T}$Q$-$I$\|$$\infty$ $\approx$ u、\allowbreak{}谱不变 & 结构 0.00e+00、\allowbreak{}正交 0.00e+00、\allowbreak{}谱 5.33e-15 \checkmark \\
A & Arnoldi 分解式误差 / 恒等式左右两端相对差 / 最大 Ritz 残差 & AV\_m $-$ V\_\{m+1\}H\_bar\_m =\allowbreak{} 0；||AV\_m y $-$ theta V\_m y|| =\allowbreak{} |h| |e\_m\textasciicircum{}T y|(\allowbreak{}同一对) & 分解式 2.22e-16；恒等式左右相对差 1.39e-16；最大 $\|$r$\|$ 4.424e-01 \checkmark \\
A & Lanczos 递推残差 / 正交性 / Ritz 值越界量 / k=n 时的谱距离 & AV\_k =\allowbreak{} V\_kT\_k + $\beta$\_\{k+1\}v\_\{k+1\}e\_k$^{\mathsf T}$；$\theta$ $\subseteq$ [$\lambda$\_min,\allowbreak{} $\lambda$\_max]；k=\allowbreak{}n 时 $\theta$ =\allowbreak{} $\lambda$ & 递推 4.44e-16、\allowbreak{}正交 1.38e-15、\allowbreak{}越界 3.55e-15、\allowbreak{}k=\allowbreak{}n 距离 9.77e-15 \checkmark \\
A & QR 迭代的谱距离 / 容许误差(6 族 $\times$ 3 位移 = 18 例) & 正交相似 $\Rightarrow$ 谱不变(\allowbreak{}病态伴随矩阵除外，见容差列) & 最坏比值 0.001；例：lap1d/\allowbreak{}wilkinson 误差 ['1.9e-14',\allowbreak{} '3.1e-15',\allowbreak{} '5.3e-15'] \checkmark \\
A & Rayleigh 商 vs 闭式 ($\lambda$$_2$$-$$\lambda$$_1$)sin$^2$$\theta$ & $\rho$(\allowbreak{}v) =\allowbreak{} $\lambda$$_1$ + (\allowbreak{}$\lambda$$_2$$-$$\lambda$$_1$)sin$^2$$\theta$(\allowbreak{}精确) & $\theta$ $\in$ [1e-1,\allowbreak{}1e-4] 上最大绝对偏差 4.44e-16 \checkmark \\
A & 自研 Jacobi SVD 与 LAPACK 的相对差 & 奇异值逐个相对误差 $\le$ 1e-12(\allowbreak{}含 $\kappa$=\allowbreak{}1e6 的对角缩放例子) & 全部奇异值最大相对差 4.95e-13；$\sigma$\_min 最大相对差 4.95e-13 \checkmark \\
A & $\sigma$\_\allowbreak{}min：单边 Jacobi vs $\surd$$\lambda$\_\allowbreak{}min(M$^{\mathsf T}$M) vs LAPACK & $\kappa$=\allowbreak{}1e12 时单边 Jacobi 仍有 \textasciitilde{}10 位；正规方程路线丢 2log$_1$$_0$$\kappa$ 位 & $\kappa$=\allowbreak{}1e+02: mine=\allowbreak{}1.000e-02,\allowbreak{} gram=\allowbreak{}1.000e-02,\allowbreak{} lapack=\allowbreak{}1.000e-02；$\kappa$=\allowbreak{}1e+04: mine=\allowbreak{}1.000e-04,\allowbreak{} gram=\allowbreak{}1.000e-04,\allowbreak{} lapack=\allowbreak{}1.000e-04；$\kappa$=\allowbreak{}1e+06: mine=\allowbreak{}1.000e-06,\allowbreak{} gram=\allowbreak{}1.000e-06,\allowbreak{} lapack=\allowbreak{}1.000e-06；$\kappa$=\allowbreak{}1e+08: mine=\allowbreak{}1.000e-08,\allowbreak{} gram=\allowbreak{}9.624e-09,\allowbreak{} lapack=\allowbreak{}1.000e-08；$\kappa$=\allowbreak{}1e+10: mine=\allowbreak{}1.000e-10,\allowbreak{} gram=\allowbreak{}0.000e+00,\allowbreak{} lapack=\allowbreak{}1.000e-10；$\kappa$=\allowbreak{}1e+12: mine=\allowbreak{}1.000e-12,\allowbreak{} gram=\allowbreak{}0.000e+00,\allowbreak{} lapack=\allowbreak{}1.000e-12 \checkmark \\
C & 幂法收敛因子(r=0.3) & 0.3 & 0.3 \checkmark \\
C & 幂法收敛因子(r=0.5) & 0.5 & 0.5 \checkmark \\
C & 幂法收敛因子(r=0.7) & 0.7 & 0.7 \checkmark \\
C & 幂法收敛因子(r=0.9) & 0.9 & 0.9 \checkmark \\
C & 幂法实测收敛因子汇总 & 每档都在 |$\lambda$$_2$/\allowbreak{}$\lambda$$_1$| 的 2\% 内 & r=\allowbreak{}0.3: 实测 0.3000；r=\allowbreak{}0.5: 实测 0.5000；r=\allowbreak{}0.7: 实测 0.7000；r=\allowbreak{}0.9: 实测 0.9000 \checkmark \\
C & 幂法 200 步 / 1e-8 的收敛阈值(二分实测) & 0.912011(\allowbreak{}=\allowbreak{} 10\textasciicircum{}(\allowbreak{}$-$8/\allowbreak{}200)) & 实测阈值 0.919433(\allowbreak{}宽 9.77e-04，9 次二分)；r=\allowbreak{}1 的旋转块 400 步后残差 1.0000，is\_complex\_pair=\allowbreak{}True \checkmark \\
C & 无位移 QR 次对角元衰减斜率 vs log|$\lambda$\_\allowbreak{}n/$\lambda$\_\allowbreak{}\{n$-$1\}| & 斜率 =\allowbreak{} log r(\allowbreak{}相对误差 $\le$ 10\%) & n=\allowbreak{}8,\allowbreak{} r=\allowbreak{}0.55: 实测 -0.5926 vs 理论 -0.5978；n=\allowbreak{}10,\allowbreak{} r=\allowbreak{}0.7: 实测 -0.3519 vs 理论 -0.3567；n=\allowbreak{}8,\allowbreak{} r=\allowbreak{}0.55: 实测 -0.5627 vs 理论 -0.5978 \checkmark \\
C & QR 迭代步数：none vs rayleigh vs wilkinson(Grcar 族) & 无位移在等模复对下不衰减；Wilkinson 位移 O(\allowbreak{}n) 步收敛 & n=\allowbreak{}8: none=\allowbreak{}600(\allowbreak{}X),\allowbreak{} rayleigh=\allowbreak{}600(\allowbreak{}X),\allowbreak{} wilkinson=\allowbreak{}27(\allowbreak{}C)；n=\allowbreak{}10: none=\allowbreak{}600(\allowbreak{}X),\allowbreak{} rayleigh=\allowbreak{}600(\allowbreak{}X),\allowbreak{} wilkinson=\allowbreak{}34(\allowbreak{}C)；n=\allowbreak{}12: none=\allowbreak{}600(\allowbreak{}X),\allowbreak{} rayleigh=\allowbreak{}600(\allowbreak{}X),\allowbreak{} wilkinson=\allowbreak{}42(\allowbreak{}C) \checkmark \\
C & 反迭代：位移 $\sigma$、实测收敛因子、理论 |$\lambda$\_\allowbreak{}next$-$$\sigma$|/|$\lambda$\_\allowbreak{}near$-$$\sigma$| & 实测因子 $\le$ 理论因子(\allowbreak{}收敛不慢于界) & $\sigma$=\allowbreak{}0.0: 45 步,\allowbreak{} 实测 0.4909,\allowbreak{} 理论 2.0000；$\sigma$=\allowbreak{}0.3: 35 步,\allowbreak{} 实测 0.3999,\allowbreak{} 理论 2.4286；$\sigma$=\allowbreak{}0.6: 25 步,\allowbreak{} 实测 0.2711,\allowbreak{} 理论 3.5000；$\sigma$=\allowbreak{}0.9: 13 步,\allowbreak{} 实测 0.0797,\allowbreak{} 理论 11.0000 \checkmark \\
C & Rayleigh 商误差斜率 / 残差斜率 & 2.0(\allowbreak{}差一个 $\theta$ 的量级) & 1.9997(\allowbreak{}残差斜率 0.5108 @ 50 点，特征值误差斜率 1.0215 @ 25 点) \checkmark \\
B & 收敛率与非正规性无关(|rate $-$ 0.5| $\le$ 0.01) & $\le$ 0.01 & $1.803\times 10^{-11}$ \checkmark \\
C & 非正规强度 $\to$ 收敛率(应不变)/ 瞬态峰谷比(应增大) & rate $\equiv$ |$\lambda$$_2$/\allowbreak{}$\lambda$$_1$| =\allowbreak{} 0.5(\allowbreak{}与 shear 无关)；瞬态随 shear 单调增 & s=\allowbreak{}0.0: rate=\allowbreak{}0.5000,\allowbreak{} transient=\allowbreak{}4.35e+12；s=\allowbreak{}0.5: rate=\allowbreak{}0.5000,\allowbreak{} transient=\allowbreak{}8.66e+12；s=\allowbreak{}1.0: rate=\allowbreak{}0.5000,\allowbreak{} transient=\allowbreak{}2.31e+13；s=\allowbreak{}2.0: rate=\allowbreak{}0.5000,\allowbreak{} transient=\allowbreak{}4.30e+13；s=\allowbreak{}5.0: rate=\allowbreak{}0.5000,\allowbreak{} transient=\allowbreak{}1.50e+14 \checkmark \\
C & Arnoldi Ritz 误差的端点行为 & m=\allowbreak{}n $\Rightarrow$ $\le$1e-12；m$\le$12 $\Rightarrow$ 误差 $\gg$ u(\allowbreak{}有界但不单调) & m=\allowbreak{}4: max=\allowbreak{}6.645e-01,\allowbreak{} mean=\allowbreak{}6.065e-01；m=\allowbreak{}6: max=\allowbreak{}3.439e-01,\allowbreak{} mean=\allowbreak{}2.994e-01；m=\allowbreak{}8: max=\allowbreak{}6.062e-01,\allowbreak{} mean=\allowbreak{}3.065e-01；m=\allowbreak{}10: max=\allowbreak{}8.422e-01,\allowbreak{} mean=\allowbreak{}3.868e-01；m=\allowbreak{}12: max=\allowbreak{}3.574e-01,\allowbreak{} mean=\allowbreak{}1.981e-01；m=\allowbreak{}16: max=\allowbreak{}9.740e-15,\allowbreak{} mean=\allowbreak{}2.666e-15 \checkmark \\
A & 对角矩阵 $\sigma$\_\allowbreak{}min vs min|z$-$$\lambda$\_\allowbreak{}i| vs LAPACK & $\sigma$\_min =\allowbreak{} dist(\allowbreak{}z,\allowbreak{} $\Lambda$)(\allowbreak{}精确，正规矩阵) & 最大相对误差(\allowbreak{}解析)1.32e-16；(\allowbreak{}LAPACK)1.67e-16 \checkmark \\
A & Jordan 块 $\sigma$\_\allowbreak{}min：与 LAPACK 的相对误差 / $\sigma$\_\allowbreak{}min $\le$ |z|\textasciicircum{}n / 地板标记数 & $\sigma$\_min =\allowbreak{} |z|\textasciicircum{}n/\allowbreak{}$\surd$(\allowbreak{}$\Sigma$|z|\textasciicircum{}\{2k\}) $\le$ |z|\textasciicircum{}n；良态点与 LAPACK 一致 & 良态点最大相对误差 5.08e-10(\allowbreak{}11 个点)；地板以下被标记 5 个 \checkmark \\
A & 正规矩阵：$\sigma$\_\allowbreak{}min 场 vs 距离场 / 伪谱掩码不一致格点数(跳过 $\varepsilon$ 边界格点) & $\sigma$\_min =\allowbreak{} dist(\allowbreak{}z,\allowbreak{}$\Lambda$)；掩码判定 等价于 dist <=\allowbreak{} epsilon(\allowbreak{}逐格点) & $\sigma$\_min 最大相对误差 2.22e-16；0 个 (\allowbreak{}格点,\allowbreak{}$\varepsilon$) 不一致 \checkmark \\
A & $\sigma$\_\allowbreak{}min 的无条件下界/上界违反量与良态点上的估计误差 & |det M|/\allowbreak{}$\|$M$\|$\_F\textasciicircum{}\{n$-$1\} $\le$ $\sigma$\_min $\le$ |t\_nn|；良态点 $\sigma$\_min 到 LAPACK 精度 & 地板以上的 18 个点：下界超出量 -3.23e-05；上界超出量 +1.11e-16；相对误差 1.26e-08 \checkmark \\
A & $\sigma$\_\allowbreak{}min 的精度与噪声地板标记($\kappa$ = 1e2 $\dots$ 1e12) & $\sigma$\_min 大于约 sqrt(\allowbreak{}u)*||M|| 时可信；以下必须置 at\_noise\_floor & $\kappa$=\allowbreak{}1e+02: mine=\allowbreak{}1.000e-02,\allowbreak{} lapack=\allowbreak{}1.000e-02,\allowbreak{} floor=\allowbreak{}0；$\kappa$=\allowbreak{}1e+04: mine=\allowbreak{}1.000e-04,\allowbreak{} lapack=\allowbreak{}1.000e-04,\allowbreak{} floor=\allowbreak{}0；$\kappa$=\allowbreak{}1e+06: mine=\allowbreak{}1.000e-06,\allowbreak{} lapack=\allowbreak{}1.000e-06,\allowbreak{} floor=\allowbreak{}0；$\kappa$=\allowbreak{}1e+08: mine=\allowbreak{}8.367e-09,\allowbreak{} lapack=\allowbreak{}1.000e-08,\allowbreak{} floor=\allowbreak{}1；$\kappa$=\allowbreak{}1e+10: mine=\allowbreak{}8.394e-09,\allowbreak{} lapack=\allowbreak{}1.000e-10,\allowbreak{} floor=\allowbreak{}1；$\kappa$=\allowbreak{}1e+12: mine=\allowbreak{}0.000e+00,\allowbreak{} lapack=\allowbreak{}1.000e-12,\allowbreak{} floor=\allowbreak{}1；良态点最大相对误差 3.00e-10；地板以下被标记 4 个；良态点被误标 0 个 \checkmark \\
A & 复数 LU 的残差与与 LAPACK 解的差 & $\|$Mx$-$b$\|$$\infty$ $\le$ u$\cdot$n$\cdot$$\|$M$\|$$\infty$$\|$x$\|$$\infty$ & 最大相对残差 2.72e-15；与 LAPACK 解的最大相对差 5.31e-16 \checkmark \\
A & Kreiss 下界 ($\rho$\_\allowbreak{}$\varepsilon$$-$1)/$\varepsilon$ 与实测峰值 sup\_\allowbreak{}k$\|$A\textasciicircum{}k$\|$ & (\allowbreak{}$\rho$\_$\varepsilon$$-$1)/\allowbreak{}$\varepsilon$ $\le$ K(\allowbreak{}A) =\allowbreak{} sup\_k$\|$A\textasciicircum{}k$\|$$_2$ & jordan: 下界 0.000e+00 $\le$ 峰值 1.000e+00；shear: 下界 2.754e+03 $\le$ 峰值 5.015e+10；grcar: 下界 6.971e+00 $\le$ 峰值 7.505e+26 \checkmark \\
A & Kreiss 常数(两条接口)与峰值位置 & K(\allowbreak{}A) =\allowbreak{} sup\_k$\|$A\textasciicircum{}k$\|$$_2$ =\allowbreak{} transient\_growth 的 peak & shear: K=\allowbreak{}3.2375e+08 @k=\allowbreak{}40,\allowbreak{} peak=\allowbreak{}3.2375e+08；shifted\_jordan: K=\allowbreak{}8.1578e+07 @k=\allowbreak{}40,\allowbreak{} peak=\allowbreak{}8.1578e+07；grcar: K=\allowbreak{}4.6293e+17 @k=\allowbreak{}40,\allowbreak{} peak=\allowbreak{}4.6293e+17；normal: K=\allowbreak{}1.0000e+00 @k=\allowbreak{}0,\allowbreak{} peak=\allowbreak{}1.0000e+00 \checkmark \\
A & GMRES 残差最大上翘量与解与 LAPACK 的差 & $\|$r\_k$\|$ 单调不增；收敛后 x 与直接解一致 & 最大上翘 -9.56e-09；解的最大相对差 5.65e-16 \checkmark \\
A & marching squares 端点上的重建值偏离 level 的最大量 & 端点插值 =\allowbreak{} level(\allowbreak{}线性插值精确) & 最大相对偏离 4.34e-16；线段数 68；零长线段 0；越界端点 0 \checkmark \\
A & 幂零 Jordan 块的 $\|$A\textasciicircum{}k$\|$ 序列 & $\|$J\_n(\allowbreak{}0)\textasciicircum{}k$\|$$_2$ =\allowbreak{} 1 (\allowbreak{}1$\le$k$\le$n$-$1)，=\allowbreak{} 0 (\allowbreak{}k$\ge$n) & 1$\le$k$\le$7 上的最大偏差 2.22e-16；k$\ge$n 的最大值 0.00e+00；$\rho$(\allowbreak{}A)=\allowbreak{}0.00e+00 \checkmark \\
C & $\rho$\_\allowbreak{}$\varepsilon$ 的幂律指数实测 vs 1/n & $\rho$\_$\varepsilon$ $\propto$ $\varepsilon$\textasciicircum{}\{1/\allowbreak{}n\}(\allowbreak{}Jordan 块) & n=\allowbreak{}6: 实测指数 0.1768 vs 1/\allowbreak{}n=\allowbreak{}0.1667(\allowbreak{}5 点)；n=\allowbreak{}10: 实测指数 0.1064 vs 1/\allowbreak{}n=\allowbreak{}0.1000(\allowbreak{}6 点) \checkmark \\
A & trusted 掩码语义 / 可信比例随网格加密的变化 & $\sigma$\_min > 100$\surd$u$\cdot$$\|$M$\|$ $\Rightarrow$ trusted；不可信格点的 $\sigma$\_min $\le$ 100$\surd$u$\cdot$$\|$M$\|$；比例保持 > 0.95 & 可信比例 0.9773(\allowbreak{}23$^2$)$\to$ 0.9737(\allowbreak{}37$^2$)；被标不可信格点中 $\sigma$\_min/\allowbreak{}floor 的最大值 4.326e-01 \checkmark \\
C & $\rho$\_\allowbreak{}$\varepsilon$ 崩塌点(扫描实测)vs 结构值 0.05\textasciicircum{}n & $\varepsilon$\_crit =\allowbreak{} 0.05\textasciicircum{}n；0.05\textasciicircum{}n $\ll$ $\surd$u 时实测提前失效(\allowbreak{}噪声决定而非结构决定) & n=\allowbreak{}4: 实测崩塌 1.000e-06 vs 结构 6.250e-06(\allowbreak{}最大 $\varepsilon$ 半径 1.2748，最小 $\varepsilon$ 半径 0.0000)；n=\allowbreak{}6: 实测崩塌 1.000e-08 vs 结构 1.563e-08(\allowbreak{}最大 $\varepsilon$ 半径 1.3176，最小 $\varepsilon$ 半径 0.0000)；n=\allowbreak{}8: 实测崩塌 1.000e-08 vs 结构 3.906e-11(\allowbreak{}最大 $\varepsilon$ 半径 1.3623，最小 $\varepsilon$ 半径 0.0000) \checkmark \\
C & 瞬态峰值对剪切强度 s 的幂律指数(有限窗口实测) & 峰值 $\propto$ s\textasciicircum{}p(\allowbreak{}p 与 n$-$1 =\allowbreak{} 7 同量级；有限 k 窗口下 p 偏小) & 实测指数 5.8849(\allowbreak{}6 点)；峰值序列 ['7.174e+04',\allowbreak{} '1.008e+06',\allowbreak{} '1.447e+07',\allowbreak{} '1.671e+08',\allowbreak{} '2.151e+09',\allowbreak{} '1.813e+10']；n$-$1 =\allowbreak{} 7 \checkmark \\
C & 「只看谱」的界的实测指数增长率 vs GMRES 残差实测下降率 & 非正规：界指数膨胀(\allowbreak{}增长率 > 0)，实测残差下降/\allowbreak{}平台(\allowbreak{}下降率 $\le$ 0) & grcar: 界增长率 +0.956/\allowbreak{}步,\allowbreak{} 残差下降率 -3.456/\allowbreak{}步,\allowbreak{} r2=\allowbreak{}8.136e-01$\to$r10=\allowbreak{}7.724e-33,\allowbreak{} bound10=\allowbreak{}3.238e+04；shear: 界增长率 +0.996/\allowbreak{}步,\allowbreak{} 残差下降率 -4.041/\allowbreak{}步,\allowbreak{} r2=\allowbreak{}2.689e-01$\to$r10=\allowbreak{}2.381e-35,\allowbreak{} bound10=\allowbreak{}4.065e+04 \checkmark \\
C & $\sigma$\_\allowbreak{}min 网格场的可复现性与相邻格点跳变 & 两次调用逐位相同(\allowbreak{}确定性)；相邻格点连续 & 逐位相同 =\allowbreak{} True；最大相邻跳变 /\allowbreak{} max $\sigma$ =\allowbreak{} 0.117 \checkmark \\
C & 非正规强度 s $\to$ GMRES 步数 / 伪谱半径 & 非正规性越强，GMRES 越慢(\allowbreak{}预解式越大) & s=\allowbreak{}0.1: 8 步,\allowbreak{} $\rho$\_$\varepsilon$=\allowbreak{}1.306,\allowbreak{} conv=\allowbreak{}1；s=\allowbreak{}0.5: 8 步,\allowbreak{} $\rho$\_$\varepsilon$=\allowbreak{}1.306,\allowbreak{} conv=\allowbreak{}1；s=\allowbreak{}1.0: 8 步,\allowbreak{} $\rho$\_$\varepsilon$=\allowbreak{}1.306,\allowbreak{} conv=\allowbreak{}1；s=\allowbreak{}2.0: 8 步,\allowbreak{} $\rho$\_$\varepsilon$=\allowbreak{}1.306,\allowbreak{} conv=\allowbreak{}1；s=\allowbreak{}4.0: 8 步,\allowbreak{} $\rho$\_$\varepsilon$=\allowbreak{}1.306,\allowbreak{} conv=\allowbreak{}1 \checkmark \\
C & $\varepsilon$\textasciicircum{}\{1/n\} 随 n 的单调性与量级 & $\varepsilon$=\allowbreak{}1e-12 时 n=\allowbreak{}32 给出 $\approx$0.42、\allowbreak{}n=\allowbreak{}64 给出 $\approx$0.65 & n=\allowbreak{}2: 0.0000；n=\allowbreak{}4: 0.0010；n=\allowbreak{}8: 0.0316；n=\allowbreak{}16: 0.1778；n=\allowbreak{}32: 0.4217 \checkmark \\
\end{longtable}
